%% =====================================================================
%%  B1-T-06-02 -- what B1 contributes to "The Two-Tellings Test"
%%  -------------------------------------------------------------------
%%  LAYER A: APPEND-ONLY. Written once, then left alone. Material found
%%  only here sits in the `unique` slot and is protected by rule R10.
%% =====================================================================
%% @kind: source
%% @id: B1-T-06-02
%% @book: B1
%% @topic: T-06-02
%% @locator: Tactic 1, pp. 14--16
%% @status: distilled
%% @unique: yes
%% @integrated: yes
%% @updated: 2026-09-19

\dossier{B1}{T-06-02}{\bookshort{B1}}{Tactic 1, pp. 14--16}

\begin{distilled}
  \item A named practitioner's method for sizing up trustworthiness: listen to the person tell the same story twice, note the details carefully the first time, and compare the two versions.
  \item The degree of consistency between versions is used as the gauge of honesty.
  \item A person who relates a story differently each time is probably an untrustworthy liar; habitual liars change stories to fit the purpose at hand and forget how they told it before.
  \item The test is presented as a practical alternative to judging character by impression.
\end{distilled}

\begin{unique}
  \item The two-tellings procedure stated as a concrete method with a defined comparison --- specifics, not impressions --- and the explanation that habitual reconstruction causes the teller to lose track of earlier versions.
\end{unique}

\begin{terms}
  \item \emph{the two-tellings test} --- this handbook's name for the procedure; the source does not label it.
\end{terms}
